% !TeX program = xelatex
\documentclass[UTF8,zihao=-4]{ctexart}

% -------------------- 页面与基础设置 --------------------
\usepackage[a4paper,margin=2.5cm]{geometry}
\linespread{1.25}\selectfont
\setlength{\parindent}{2em}

% -------------------- 数学与排版 --------------------
\usepackage{amsmath,amssymb,amsfonts}
\usepackage{enumitem}
\usepackage{graphicx}
\usepackage{float}
\usepackage[table]{xcolor}
\usepackage{booktabs}
\usepackage{tabularx}
\usepackage{array}
\usepackage{caption}
\usepackage{subcaption}
\usepackage{listings}
\usepackage{fancyhdr}
\usepackage{url}
\usepackage[hidelinks]{hyperref}

% -------------------- 基本信息宏 --------------------
\newcommand{\ReportCourse}{2026 机器学习}
\newcommand{\ReportNo}{1}
\newcommand{\ReportName}{徐厚泽}
\newcommand{\ReportStudentID}{26113050327}

% -------------------- 页眉页脚 --------------------
\pagestyle{fancy}
\fancyhf{}
\fancyfoot[C]{\thepage}
\renewcommand{\headrulewidth}{0pt}

% -------------------- 代码样式 --------------------
\lstdefinestyle{PythonCode}{
  language=Python,
  numbers=left,
  numberstyle=\tiny\color{gray},
  xleftmargin=2.2em,
  basicstyle=\ttfamily\small,
  keywordstyle=\color{blue!70!black}\bfseries,
  stringstyle=\color{green!40!black},
  commentstyle=\color{gray!70!black},
  breaklines=true,
  breakatwhitespace=false,
  columns=flexible,
  frame=single,
  rulecolor=\color{black!20},
  backgroundcolor=\color{black!3}
}
\lstset{style=PythonCode}

% -------------------- 图表与列表样式 --------------------
\captionsetup{font=small,labelfont=bf,labelsep=quad}
\setlist{itemsep=0.25em,topsep=0.25em,parsep=0pt,partopsep=0pt}

% -------------------- 常用自定义命令 --------------------
\newcommand{\mlExperiment}[2]{\subsection{实验#1：#2}}
\newcommand{\mlConclusion}[1]{\par\noindent\textbf{结论#1：}}

\begin{document}

% -------------------- 封面 --------------------
\begin{titlepage}
  \centering
  \vspace*{1.2cm}
  {\heiti\zihao{1}\ReportCourse\par}
  \vspace{1.5em}
  {\heiti\zihao{1}课程报告\ReportNo\par}
  \vfill
  {\kaishu\zihao{-3}姓名：\ReportName\par}
  \vspace{0.8em}
  {\kaishu\zihao{-3}学号：\ReportStudentID\par}
  \vfill
\end{titlepage}

% -------------------- 正文 --------------------
\section{任务描述}
简要描述本次作业的任务目标、要解决的问题、数据来源以及评价指标。

\section{数据描述}
介绍数据集的来源、样本规模、特征含义、训练集/测试集划分方式等信息。

\section{数据预处理}
说明数据格式转换、归一化/标准化、采样、降维、数据增强等预处理操作，并给出关键代码。

\begin{lstlisting}
import numpy as np

# 示例：读入数据并做预处理
x_train = np.load("x_train.npy")

x_train = x_train.reshape(len(x_train), -1).astype("float32") / 255.
\end{lstlisting}

\section{算法介绍}
\subsection{算法原理}
说明所实现或对比的算法的基本思想、适用场景和关键步骤。

\subsection{算法实现}
说明算法的实现流程，并给出伪代码、核心代码或模块划分。

\begin{lstlisting}
def predict(self, x_test):
    # 示例：实现核心预测流程
    results = []
    return np.array(results)
\end{lstlisting}

\section{实验结果及分析}

\mlExperiment{5.1}{baseline}
说明 baseline 实验设置，包括训练样本数量、测试样本数量、关键超参数和评价方式。

\begin{figure}[H]
  \centering
  \fbox{\parbox[c][4cm][c]{0.68\textwidth}{\centering 请在此放置 baseline 实验图}}
  \caption{baseline 实验结果}
  \label{fig:baseline}
\end{figure}

\begin{table}[H]
  \centering
  \caption{baseline 实验结果对比}
  \label{tab:baseline}
  \begin{tabular}{lcc}
    \toprule
    方法 & 准确率 & 运行时间/s \\
    \midrule
    baseline & 0.0000 & 0.000 \\
    \bottomrule
  \end{tabular}
\end{table}

\mlConclusion{1}
根据实验结果总结第一条结论，例如方法是否有效、瓶颈在哪里、是否可以进一步优化等。

\mlExperiment{5.2}{改进实验}
说明改进实验的动机、设置和实现方式。

\mlConclusion{2}
根据实验结果总结第二条结论。

\section{总结}
总结本次作业完成的工作、得到的主要结论、遇到的问题以及后续可以改进的方向。

% 如果需要参考文献，可以在此添加 references.bib，并取消以下注释：
% \bibliographystyle{plain}
% \bibliography{references}

\end{document}
